\subsection{REGULAR SPACES}

PROPOSITION 11. The following properties of a topological space X are equivalent:

(O \(_{III}\) ) The set of closed neighbourhoods of any point of X is a fundamental system of neighbourhoods of the point.

\((O_{\mathrm{III}}^{\prime})\) Given any closed subset \(\mathbf{F}\) of \(\mathbf{X}\) and any point \(x \notin \mathbf{F}\) there is a neighbourhood of \(x\) and a neighbourhood of \(\mathbf{F}\) which do not intersect.

\((\mathrm{O}_{\mathrm{III}}) \Longrightarrow (\mathrm{O}_{\mathrm{III}}^{\prime})\) : If \(\mathbf{F}\) is closed and \(x \notin \mathbf{F}\) , then there is a closed neighbourhood \(\mathbf{V}\) of \(x\) contained in the neighbourhood \(\complement \mathbf{F}\) of \(x\) ; \(\mathbf{V}\) and \(\complement \mathbf{V}\) are neighbourhoods of \(x\) and \(\mathbf{F}\) respectively, and have no point in common.

\((\mathrm{O}_{\mathrm{III}}^{\prime}) \Longrightarrow (\mathrm{O}_{\mathrm{III}})\) : If \(W\) is an open neighbourhood of \(x \in X\) , then there is a neighbourhood \(U\) of \(x\) and a neighbourhood \(V\) of \(\complement W\) which are disjoint, and therefore \(\overline{U} \subset W\) .

DEFINITION 2. A topological space is said to be regular if it is Hausdorff and satisfies axiom (O \(_{III}\) ); its topology is then said to be regular.

A discrete space is regular. * We shall see in § 9 that every locally compact space (in particular the real line R) is regular. *

PROPOSITION 12. Every subspace of a regular space is regular.

Let A be a subspace of a regular space X. Since X is Hausdorff, so is A (no. 2); on the other hand, every neighbourhood of a point \(x \in A\) with respect to A is of the form \(V \cap A\) , where V is a neighbourhood of x in X. Since X is regular there is a neighbourhood W of x in X which is closed in X and contained in V; \(W \cap A\) is then a neighbourhood of x in A, closed in A and contained in \(V \cap A\) . Hence the result. Conversely:

PROPOSITION 13. If every point x of a topological space X has a closed neighbourhood which is a regular subspace of X, then X is regular.

X is Hausdorff by Proposition 6 of no. 2. Let x be any point of X and let V be a closed regular neighbourhood of x. If U is any neighbourhood of x contained in V, then U is a neighbourhood of x relative to V; hence by hypothesis there is a neighbourhood W of x in V which is closed in V and contained in U. But W is a neighbourhood of x in X since V is a neighbourhood of x in X, and W is closed in X since V is closed in X.

Remarks. 1) There are examples of non-Hausdorff spaces in which every point has a regular neighbourhood (Exercise 7).

\begin{enumerate}
\def\labelenumi{\arabic{enumi})}
\setcounter{enumi}{1}
\item
  There are spaces which are Hausdorff but not regular (Exercise 20).
\item
  A topology which is finer than a regular topology need not be regular (Exercise 20).
\end{enumerate}

